% ==========================================
% CHAPTER 3
% ==========================================
\chapter{Libraries and Basic Script Writing}

\section{Introduction}
We saw earlier that one can use the iPython interpreter to do basic math, and that there were various data types that come ``preinstalled'' within Python (like lists, strings, integers, etc). However, once a code requires more sophisticated analytical tools (especially for astronomical processes), it becomes apparent that the vanilla iPython functions are not sufficient. Luckily, there are hundreds of functions that have been written to accomplish these tasks, most of which are organized into what are called libraries.

\textbf{Definition 3.0.1} A library is a maintained collection of functions which can be installed and used.

Most Python distributions come with a lot of these libraries included, and installing new libraries is generally straightforward. There are 4 key libraries that we will be discussing in detail in this text: \texttt{numpy}, \texttt{matplotlib}, and (sometimes) \texttt{scipy} and \texttt{astropy}. Numpy is an extremely versatile library of functions to do the things Python can't. Matplotlib, meanwhile, is a library with functions dedicated to plotting data and making graphs. Astropy is a library with functions specifically for astronomical applications: we will be using it to import FITS images (images taken by telescopes), among other things. Scipy is a library that contains special use functions that are often used in science.

\section{Installing Libraries}
Most scientific distributions of Python (like Anaconda) come with important packages like numpy preinstalled. However, for most smaller packages, like astropy, or pyfits, you will likely have to install them yourself. The easiest way, when available, is to use pip:

\begin{lstlisting}
>> pip install packagename
\end{lstlisting}

\section{Importing Libraries}
Because these libraries are not automatically loaded up when Python runs, we have to import them:

\begin{lstlisting}
[IN]: import numpy
[IN]: import matplotlib.pyplot
[IN]: import astropy.io.fits
\end{lstlisting}

Clearly, writing out \texttt{numpy} all over your code would take forever. Luckily, python allows us to import the libraries and name them whatever we want for the purposes of our code. Two standard choices are:

\begin{lstlisting}
[IN]: import numpy as np
[IN]: import matplotlib.pyplot as plt
\end{lstlisting}

\section{Writing A Program}
Thus far, we have been working entirely in the iPython interpreter. While this is a quick and easy way to practice with Python, it is unsuitable for the majority of things that you might want to accomplish. Thus, most of the times we write what are known as scripts, or programs.

\textbf{Definition 3.3.1} A program is a self-contained list of commands that are stored in a file that can be read by Python. Essentially, it is a text file, with each line being the exact syntax you would have typed into the terminal.

\begin{lstlisting}
import numpy as np
import matplotlib.pyplot as plt

x = np.arange(100)
y = x**2 + np.sin(3*x)

plt.plot(x,y)
plt.show()
\end{lstlisting}

\textbf{Exercise 3.1 Loading 1-dimensional data from a file \& plotting} \\
The first step needed to work with any astronomical data is, of course, to load it into your code. For that, we utilize a numpy function called \texttt{loadtxt}. Let's say we have a file on our computer called \texttt{spectrum.txt} that contains two columns, wavelength and flux:

\begin{lstlisting}
data = np.loadtxt('spectrum.txt')
data = np.transpose(data)
fluxes, wavelength = data[0], data[1]

plt.plot(wavelength, fluxes)
plt.show()
\end{lstlisting}

\section{Working with Arrays}
\subsection{Creating a Numpy Array}
Here's a bunch of ways to initialize a basic numpy array:

\begin{lstlisting}
empty = np.array([])
zeros = np.zeros(len_desired) # creates an array of zeros
ones = np.ones(len_desired) # creates an array of 1's
twos = np.ones(len_desired)*2 # creates an array of 2's
count = np.arange(start,stop,step) # creates an array of integers
resolution = np.linspace(start,stop,num) # creates an array of floats
\end{lstlisting}

\subsection{Basic Array Manipulation}
If you need to change a value in an array, the syntax is identical to before, simply set \texttt{arrayname[index] = new\_value} to change it. To delete values from an array, you can use:

\begin{lstlisting}
arrayname = np.delete(arrayname, indices)
\end{lstlisting}